\newcommand{\serverglyph}[1]{%
\begin{scope}[shift={#1}]
  \draw[fill=gray!20, rounded corners=1pt] (0,0) rectangle (0.36,0.62);
  \draw[fill=gray!35] (0.36,0) -- (0.52,0.10) -- (0.52,0.72) -- (0.36,0.62) -- cycle;
  \draw (0.07,0.46) -- (0.28,0.46);
  \draw (0.07,0.34) -- (0.28,0.34);
  \draw (0.07,0.22) -- (0.20,0.22);
  \fill (0.09,0.09) circle (0.025);
\end{scope}}

\newcommand{\dbglyph}[1]{%
\begin{scope}[shift={#1}]
  \draw[fill=gray!20] (0,0) ellipse (0.34 and 0.11);
  \draw[fill=gray!12] (-0.34,0) -- (-0.34,-0.55) arc (180:360:0.34 and 0.11) -- (0.34,0);
  \draw (0,-0.18) ellipse (0.34 and 0.11);
  \draw (0,-0.37) ellipse (0.34 and 0.11);
\end{scope}}

\newcommand{\personglyph}[1]{%
\begin{scope}[shift={#1}]
  \draw[fill=gray!10] (0,0) circle (0.34);
  \fill[gray!55] (0,0.12) circle (0.11);
  \draw[fill=gray!45] (-0.18,-0.18) arc (180:0:0.18 and 0.17) -- (0.18,-0.28) -- (-0.18,-0.28) -- cycle;
\end{scope}}

\newcommand{\thesisprismatikz}{%
\begin{figure}[H]
\centering
\resizebox{0.74\textwidth}{!}{%
\begin{tikzpicture}[
  box/.style={rectangle, draw, very thick, minimum width=9.2cm, minimum height=1.55cm, align=center, font=\large},
  arrow/.style={-{Latex[length=4mm]}, line width=1.2pt},
  node distance=0.72cm
]
\node[font=\bfseries\LARGE] (title) {PRISMA Flow Diagram};
\node[box, below=0.85cm of title] (id) {Identification\\{\normalsize 2581 records retrieved}};
\node[box, below=of id] (screen) {Screening\\{\normalsize 5 duplicates removed}\\{\normalsize 2576 records screened}};
\node[box, below=of screen] (elig) {Eligibility\\{\normalsize 2559 excluded}\\{\normalsize 17 full-text assessed}};
\node[box, below=of elig] (inc) {Included\\{\normalsize 6 studies in final review}};
\draw[arrow] (id) -- (screen);
\draw[arrow] (screen) -- (elig);
\draw[arrow] (elig) -- (inc);
\end{tikzpicture}}
\caption{PRISMA flow diagram.}
\end{figure}}

\newcommand{\thesislayerstikz}{%
\begin{figure}[H]
\centering
\resizebox{0.92\textwidth}{!}{%
\begin{tikzpicture}[
  baseLayer/.style={rectangle, rounded corners=6pt, very thick, minimum width=10.5cm, minimum height=1.55cm, text width=10cm, align=center},
  appLayer/.style={baseLayer, draw=orange!70!black, fill=orange!18},
  serviceLayer/.style={baseLayer, draw=blue!70!black, fill=blue!18},
  overlayLayer/.style={baseLayer, draw=green!70!black, fill=green!18},
  physicalLayer/.style={baseLayer, draw=gray!70!black, fill=gray!18},
  label/.style={font=\bfseries, align=right, text width=2.35cm},
  node distance=0.42cm
]
\node[appLayer] (app) {\textbf{Application Layer}\\[2pt]Decentralized Applications \quad Social Network\\File Sharing \quad DApps};
\node[serviceLayer, below=of app] (svc) {\textbf{Service Layer}\\[2pt]Secure Messaging \quad Data Access Services \quad Session Management};
\node[overlayLayer, below=of svc] (ovl) {\textbf{Overlay Network Layer}\\[2pt]Peer Discovery \quad Routing \& Anonymity \quad Distributed Storage};
\node[physicalLayer, below=of ovl] (phy) {\textbf{Physical \& Classical Network Layer}\\[2pt]IP / UDP / TCP / QUIC \quad NAT Traversal \quad STUN / TURN / Relay};
\node[label, left=0.8cm of app] {Application-Level};
\node[label, left=0.8cm of svc] {Reusable Services};
\node[label, left=0.8cm of ovl] {Overlay Coordination};
\node[label, left=0.8cm of phy] {Physical Connectivity};
\end{tikzpicture}}
\caption{Layered architecture.}
\label{fig:layered-architecture}
\end{figure}}

\newcommand{\thesisconnectiontikz}{%
\begin{figure}[H]
\centering
\resizebox{0.98\textwidth}{!}{%
\begin{tikzpicture}[
  panel/.style={rectangle, draw, minimum width=2.75cm, minimum height=1.45cm, align=center},
  main/.style={rectangle, draw, rounded corners=5pt, very thick, minimum width=4.25cm, minimum height=0.9cm, align=center, font=\bfseries},
  small/.style={rectangle, draw, minimum width=2.35cm, minimum height=0.62cm, align=center, font=\scriptsize},
  arrow/.style={-{Latex[length=3mm]}, thick},
  dashedlink/.style={dashed, gray!70, thick}
]
\node[font=\bfseries\Large] at (0,4.1) {Connection Management};
\draw (-6.4,3.72) -- (3.35,3.72);
\node[panel] (phys) at (-5.0,2.25) {\textbf{Physical Peers}\\[0.75cm]};
\serverglyph{(-5.45,1.77)} \serverglyph{(-4.65,1.77)}
\node[panel, text width=2.45cm] (meta) at (-5.0,-0.15) {\textbf{Peer Metadata}\\[3pt]\scriptsize endpoints\\reachability\\transport};
\node[main] (cm) at (0,2.25) {Connection Management};
\node[panel, text width=2.45cm] (endpoint) at (5.05,3.0) {\textbf{Endpoint Options}\\[3pt]IPv4\\IPv6\\Relay};
\node[panel, text width=2.45cm] (status) at (5.05,-0.15) {\textbf{Status Check}\\[3pt]Valid\\Degraded\\Failed};
\node[align=left, text width=5.85cm, font=\small] (opts) at (0,0.0) {$\bullet$ Direct Connection\\$\bullet$ Hole Punching / NAT Traversal\\$\bullet$ Relay-Assisted Establishment\\$\bullet$ Continuous Relay Forwarding};
\node[font=\bfseries] (auto) at (0,-2.22) {Automatic Selection};
\node[small] (ka) at (-3.75,-3.55) {Keep-Alive};
\node[small] (fd) at (-1.25,-3.55) {Failure Detection};
\node[small] (ru) at (1.25,-3.55) {Retry \& Update};
\node[small] (rf) at (3.75,-3.55) {Relay Forwarding};
\draw[arrow] (phys.east) -- (cm.west);
\draw[arrow] (cm.east) -- (endpoint.west);
\draw[dashedlink] (phys) -- (meta);
\draw[dashedlink] (cm.east) -- (status.west);
\draw[arrow] (opts.south) -- (auto.north);
\draw[arrow] (auto.south) -- ++(0,-0.45) -| (ka.north);
\draw[arrow] (auto.south) -- ++(0,-0.45) -| (fd.north);
\draw[arrow] (auto.south) -- ++(0,-0.45) -| (ru.north);
\draw[arrow] (auto.south) -- ++(0,-0.45) -| (rf.north);
\end{tikzpicture}}
\caption{Connection Management diagram.}
\end{figure}}

\newcommand{\thesisdrttikz}{%
\begin{figure}[H]
\centering
\resizebox{0.98\textwidth}{!}{%
\begin{tikzpicture}[
  vnode/.style={rectangle, draw=blue!55!black, rounded corners=5pt, fill=blue!8, minimum width=1.9cm, minimum height=0.72cm, align=center},
  pgroup/.style={rectangle, draw=gray!75!black, dashed, rounded corners=8pt, fill=gray!4, minimum width=3.15cm, minimum height=1.18cm, align=center},
  pn/.style={rectangle, draw=gray!70!black, rounded corners=3pt, fill=gray!16, minimum width=1.18cm, minimum height=0.38cm, align=center, font=\scriptsize},
  cell/.style={font=\scriptsize, align=center},
  arrow/.style={-{Latex[length=3mm]}, thick},
  dashedarrow/.style={-{Latex[length=3mm]}, dashed, thick}
]
\node[font=\bfseries\Large] at (0,5.35) {Distributed Routing Table (DRT)};
\node[text width=10.2cm, align=center, font=\small] at (0,4.55) {The DRT indirectly maps a virtual node identifier to physical entry points authorized to receive its traffic.};
\node[align=center] at (-5.85,3.08) {\textbf{Virtual Nodes}\\(logical identities)};
\node[vnode] (v1) at (-5.85,2.05) {$H(pk_1)$};
\node[vnode] (v2) at (-5.85,0.60) {$H(pk_2)$};
\node[vnode] (vn) at (-5.85,-1.10) {$H(pk_n)$};
\node[draw, rounded corners=8pt, fill=gray!5, minimum width=5.75cm, minimum height=3.82cm] (table) at (0,0.45) {};
\node[font=\bfseries] at (0,2.58) {DRT Records};
\draw (-2.55,1.68) rectangle (2.55,-1.02);
\draw (-0.35,1.68) -- (-0.35,-1.02);
\draw (-2.55,1.0) -- (2.55,1.0);
\draw (-2.55,0.32) -- (2.55,0.32);
\draw (-2.55,-0.36) -- (2.55,-0.36);
\node[cell, font=\scriptsize\bfseries] at (-1.45,1.34) {Virtual ID};
\node[cell, font=\scriptsize\bfseries] at (1.10,1.34) {Entry point IDs};
\node[cell] at (-1.45,0.66) {$H(pk_1)$};
\node[cell] at (1.10,0.66) {$\{pn_a,\ pn_b\}$};
\node[cell] at (-1.45,-0.02) {$H(pk_2)$};
\node[cell] at (1.10,-0.02) {$\{pn_c,\ pn_d\}$};
\node[cell] at (-1.45,-0.70) {$H(pk_n)$};
\node[cell] at (1.10,-0.70) {$\{pn_x,\ ...\}$};
\node[align=center] at (5.75,3.08) {\textbf{Physical Nodes}\\(entry points)};
\node[pgroup] (p1) at (5.75,1.90) {};
\node[pn] at (5.10,2.08) {$pn_a$};
\node[pn] at (6.40,2.08) {$pn_b$};
\node[cell] at (5.75,1.60) {authorized for $H(pk_1)$};
\node[pgroup] (p2) at (5.75,0.40) {};
\node[pn] at (5.10,0.58) {$pn_c$};
\node[pn] at (6.40,0.58) {$pn_d$};
\node[cell] at (5.75,0.10) {authorized for $H(pk_2)$};
\node[pgroup] (p3) at (5.75,-1.10) {};
\node[pn] at (5.10,-0.92) {$pn_x$};
\node[pn] at (6.40,-0.92) {$...$};
\node[cell] at (5.75,-1.40) {authorized for $H(pk_n)$};
\draw[dashedarrow] (v1.east) -- (table.west |- v1);
\draw[dashedarrow] (v2.east) -- (table.west |- v2);
\draw[dashedarrow] (vn.east) -- (table.west |- vn);
\draw[arrow] (table.east |- p1) -- (p1.west);
\draw[arrow] (table.east |- p2) -- (p2.west);
\draw[arrow] (table.east |- p3) -- (p3.west);
\node[draw, rounded corners=5pt, text width=10.2cm, align=left, fill=blue!4, font=\small] at (0,-3.35) {$\bullet$ The DRT does not store complete routes.\\$\bullet$ It provides authorized physical entry points.\\$\bullet$ Virtual-to-physical mapping remains indirect.};
\end{tikzpicture}}
\caption{Distributed Routing Table.}
\end{figure}}

\newcommand{\thesishopbyhoptikz}{%
\begin{figure}[H]
\centering
\resizebox{0.98\textwidth}{!}{%
\begin{tikzpicture}[
  hop/.style={rectangle, draw=gray!75!black, rounded corners=5pt, fill=gray!10, minimum width=1.45cm, minimum height=0.74cm, align=center, font=\scriptsize},
  vnode/.style={rectangle, draw=blue!60!black, rounded corners=6pt, fill=blue!8, minimum width=1.7cm, minimum height=0.78cm, align=center, font=\scriptsize},
  finalnode/.style={rectangle, draw=red!65!black, rounded corners=6pt, fill=red!6, minimum width=1.65cm, minimum height=0.78cm, align=center, font=\scriptsize},
  route/.style={-{Latex[length=3mm]}, very thick, blue!65!black},
  alt/.style={-{Latex[length=2.6mm]}, dashed, thick, gray!70},
  note/.style={font=\scriptsize, align=center, fill=white, inner sep=1.5pt}
]
\node[font=\bfseries\Large] at (0,5.45) {TTL-Based Hop-by-Hop Route Example};
\node[text width=10.1cm, align=center, font=\small] at (0,4.45) {A route is built incrementally: each physical node stores only its local mapping and chooses the next hop while TTL remains valid.};
\node[vnode] (vn) at (-5.85,2.45) {VN\\request};
\node[hop] (entry) at (-3.55,2.45) {Entry PN\\TTL 4};
\node[hop] (hop1) at (-1.25,2.45) {Hop 1\\TTL 3};
\node[hop] (hop2) at (1.05,2.45) {Hop 2\\TTL 2};
\node[hop] (hop3) at (3.35,2.45) {Hop 3\\TTL 1};
\node[finalnode] (final) at (5.65,2.45) {Final PN\\TTL 0};
\draw[route] (vn.east) -- (entry.west);
\draw[route] (entry.east) -- (hop1.west);
\draw[route] (hop1.east) -- (hop2.west);
\draw[route] (hop2.east) -- (hop3.west);
\draw[route] (hop3.east) -- (final.west);
\node[note] at (-4.70,2.90) {entry};
\node[note] at (-2.40,2.90) {random next};
\node[note] at (-0.10,2.90) {random next};
\node[note] at (2.20,2.90) {random next};
\node[note, text=red!70!black] at (4.50,2.90) {TTL ends};
\node[draw=blue!60!black, rounded corners=5pt, fill=blue!5, align=center, font=\scriptsize, text width=3.0cm] at (-3.45,1.05) {Each hop saves only\\previous $\rightarrow$ next};
\node[draw=blue!60!black, rounded corners=5pt, fill=blue!5, align=center, font=\scriptsize, text width=3.0cm] at (0,1.05) {TTL decreases\\at every hop};
\node[draw=red!65!black, rounded corners=5pt, fill=red!5, align=center, font=\scriptsize, text width=3.0cm] at (3.45,1.05) {Construction finishes\\at the final PN};
\draw[alt] (-3.45,0.25) .. controls (-1.25,-0.42) and (1.25,-0.42) .. (3.45,0.25);
\node[note] at (0,-0.40) {another random walk can build a different path};
\node[draw, rounded corners=5pt, fill=blue!4, text width=10.2cm, align=left, font=\small] at (0,-2.25) {$\bullet$ Each hop chooses one next physical node.\\$\bullet$ TTL decreases at every forwarding step.\\$\bullet$ The route emerges during forwarding.};
\end{tikzpicture}}
\caption{Example of hop-by-hop route construction with a TTL-based strategy.}
\label{fig:ttl-route-strategy-example}
\end{figure}}

\newcommand{\thesisroutestrategyconstructiontikz}{%
\begin{figure}[H]
\centering
\resizebox{0.98\textwidth}{!}{%
\begin{tikzpicture}[
  block/.style={rectangle, draw=gray!75!black, rounded corners=6pt, fill=gray!8, minimum width=2.05cm, minimum height=0.72cm, align=center, font=\scriptsize},
  protocol/.style={rectangle, draw=blue!65!black, rounded corners=7pt, fill=blue!8, minimum width=2.75cm, minimum height=0.88cm, align=center, font=\scriptsize},
  strategy/.style={rectangle, draw=orange!75!black, rounded corners=7pt, fill=orange!10, minimum width=2.75cm, minimum height=0.88cm, align=center, font=\scriptsize},
  hop/.style={circle, draw=gray!70!black, fill=gray!8, minimum size=0.88cm, align=center, font=\scriptsize},
  finalnode/.style={circle, draw=red!65!black, fill=red!8, minimum size=0.98cm, align=center, font=\scriptsize},
  record/.style={rectangle, draw=purple!65!black, rounded corners=5pt, fill=purple!7, minimum width=2.45cm, minimum height=0.68cm, align=center, font=\scriptsize},
  arrow/.style={-{Latex[length=2.8mm]}, thick, blue!65!black},
  control/.style={-{Latex[length=2.5mm]}, thick, orange!80!black},
  dashedarrow/.style={-{Latex[length=2.5mm]}, dashed, thick, gray!70},
  note/.style={font=\scriptsize, align=center, fill=white, inner sep=1.2pt}
]
\node[font=\bfseries\Large] at (0,5.70) {Strategy-Based Route Construction};
\node[text width=11.6cm, align=center, font=\small] at (0,4.72) {The route protocol keeps the lifecycle stable. The active route strategy chooses each next physical hop and defines when construction should finish.};

\node[block, fill=green!8, draw=green!45!black] (vn) at (-5.35,3.10) {Virtual node\\requests route};
\node[protocol] (protocol) at (-2.55,3.10) {Route protocol\\lifecycle};
\node[strategy] (strategy) at (0.45,3.10) {Route strategy\\policy};
\node[block, fill=blue!6, draw=blue!55!black] (mapping) at (3.25,3.10) {Local mappings\\per hop};
\node[record, minimum width=2.15cm] (drt) at (6.05,3.10) {DRT entry\\published};

\draw[arrow] (vn.east) -- (protocol.west);
\draw[arrow] (protocol.east) -- (strategy.west);
\draw[arrow] (strategy.east) -- (mapping.west);
\draw[arrow] (mapping.east) -- (drt.west);

\node[hop] (h1) at (-4.45,1.30) {PN};
\node[hop] (h2) at (-2.20,1.30) {PN};
\node[hop] (h3) at (0.05,1.30) {PN};
\node[hop] (h4) at (2.30,1.30) {PN};
\node[finalnode] (final) at (4.55,1.30) {Final\\PN};

\draw[arrow] (h1.east) -- node[note, above=3pt, font=\tiny] {\texttt{path\_id}} (h2.west);
\draw[arrow] (h2.east) -- node[note, above=3pt, font=\tiny] {\texttt{path\_id}} (h3.west);
\draw[arrow] (h3.east) -- node[note, above=3pt, font=\tiny] {\texttt{path\_id}} (h4.west);
\draw[arrow] (h4.east) -- node[note, above=3pt, font=\tiny] {\texttt{path\_id}} (final.west);

\draw[control] (strategy.south) -- ++(0,-0.55) -| node[pos=0.28, right, note] {next-hop choice} (h3.north);
\draw[dashedarrow] (final.north) -- ++(0,0.55) -| (drt.south);

\node[record, text width=2.40cm] (candidates) at (-4.00,-0.35) {Candidate peers\\and local state};
\node[record, text width=2.65cm] (objectives) at (-0.55,-0.35) {Optimization goal\\privacy / latency\\bandwidth / reliability};
\node[record, text width=2.40cm] (termination) at (2.90,-0.35) {Termination rule\\strategy-specific};

\node[draw, rounded corners=5pt, fill=blue!4, text width=11.0cm, align=left, font=\small] at (0,-2.62) {$\bullet$ The protocol defines invariant route lifecycle messages and publication rules.\\$\bullet$ The strategy defines hop selection and stopping conditions.\\$\bullet$ Different strategies can optimize privacy, latency, bandwidth, reliability, or diversity without changing route execution.};
\end{tikzpicture}}
\caption{Generic route construction through a route strategy.}
\end{figure}}

\newcommand{\thesispeerdiscoverytikz}{%
\begin{figure}[H]
\centering
\resizebox{0.96\textwidth}{!}{%
\begin{tikzpicture}[
  arrow/.style={-{Latex[length=2.8mm]}, thick, blue!65!black},
  dashedarrow/.style={-{Latex[length=2.5mm]}, dashed, thick, gray!70},
  seed/.style={rectangle, draw=green!45!black, rounded corners=4pt, fill=green!8, minimum width=2.38cm, minimum height=0.58cm, align=center, font=\scriptsize},
  public/.style={circle, draw=blue!60!black, fill=blue!8, minimum size=0.9cm, align=center, font=\scriptsize},
  private/.style={circle, draw=red!60!black, fill=red!7, minimum size=0.9cm, align=center, font=\scriptsize},
  box/.style={rectangle, draw=gray!75!black, rounded corners=5pt, fill=gray!8, minimum width=2.35cm, minimum height=0.78cm, align=center, font=\scriptsize},
  relay/.style={diamond, draw=orange!70!black, fill=orange!10, aspect=1.6, inner sep=1.5pt, minimum width=1.55cm, minimum height=1.0cm, align=center, font=\scriptsize},
  meta/.style={rectangle, draw=purple!55!black, rounded corners=5pt, fill=purple!7, minimum width=2.85cm, minimum height=0.78cm, text width=2.65cm, align=center, font=\scriptsize},
  stage/.style={font=\bfseries\scriptsize, align=center, text width=2.55cm}
]
\node[font=\bfseries\Large] at (0,5.85) {Peer Discovery Process};
\node[text width=13.0cm, align=center, font=\small] at (0,4.82) {Public peers bootstrap the network, learn about private peers, and share reachability metadata for relay or hole-punch attempts.};

\node[stage] at (-5.55,3.78) {1. Bootstrap};
\node[seed] (seed1) at (-5.55,2.92) {DNS seeds};
\node[seed] (seed2) at (-5.55,2.05) {Hardcoded\\public nodes};

\node[stage] at (-2.55,3.78) {2. Public peer view};
\node[public] (pubA) at (-3.25,2.68) {P1};
\node[public] (pubB) at (-1.85,2.68) {P2};
\node[public] (pubC) at (-2.55,1.65) {P3};
\node[box, fill=blue!6, draw=blue!50!black, minimum width=2.75cm] (publicView) at (-2.55,0.35) {Known public\\peer exchange};
\draw[dashed, gray!65] (pubA) -- (pubB) -- (pubC) -- (pubA);

\node[stage] at (0.90,3.78) {3. Private peers};
\node[box, draw=red!55!black, fill=red!4, minimum width=3.05cm,
      minimum height=1.85cm] (nat) at (0.90,2.05) {};
\node[private] (privA) at (0.20,2.30) {N1};
\node[private] (privB) at (1.60,2.30) {N2};
\node[font=\scriptsize] at (0.90,1.43) {Behind NAT / Firewall};

\node[stage] at (4.85,3.78) {4. Connect};
\node[meta] (meta) at (4.85,2.45) {Published metadata\\reachable via relay\\or hole punching};
\node[relay] (relay) at (3.82,0.62) {Relay};
\node[box, fill=orange!7, draw=orange!70!black] (hole) at (5.88,0.62) {Hole punching\\attempt};
\node[public] (newPeer) at (4.85,-0.72) {P4};

\draw[arrow] (seed1.east) -- (pubA.west);
\draw[arrow] (seed2.east) -- (pubC.west);
\draw[arrow] (pubA.south) -- (publicView.north west);
\draw[arrow] (pubB.south) -- (publicView.north east);
\draw[arrow] (pubC.south) -- (publicView.north);

\draw[arrow, shorten >=4pt, shorten <=4pt] (nat.south west) -- (publicView.east);
\draw[arrow, shorten >=4pt, shorten <=4pt] (nat.east) -- (meta.west);
\coordinate (connectBranch) at (4.85,1.25);
\draw[arrow] (meta.south) -- (connectBranch) -| (relay.north);
\draw[arrow] (meta.south) -- (connectBranch) -| (hole.north);
\draw[arrow] (relay.south) -- (newPeer.north west);
\draw[dashedarrow] (hole.south) -- (newPeer.north east);

\node[draw, rounded corners=5pt, fill=blue!4, text width=11.6cm, align=left, font=\small] at (0,-2.45) {$\bullet$ New peers first contact DNS seeds or hardcoded public nodes.\\$\bullet$ Private peers connect to public peers and publish metadata.\\$\bullet$ Other peers use that metadata for relay routing or NAT hole punching.};
\end{tikzpicture}}
\caption{Peer discovery process.}
\end{figure}}

\newcommand{\thesisddttikz}{%
\begin{figure}[H]
\centering
\resizebox{0.98\textwidth}{!}{%
\begin{tikzpicture}[
  arrow/.style={-{Latex[length=3mm]}, thick, blue!60!black},
  dashedarrow/.style={-{Latex[length=2.6mm]}, dashed, thick, gray!70},
  vnode/.style={circle, draw=gray!75!black, fill=gray!10, minimum size=0.95cm, align=center, font=\scriptsize},
  tablebox/.style={draw=gray!80!black, rounded corners=6pt, fill=gray!4},
  card/.style={draw=blue!55!black, rounded corners=5pt, fill=blue!5, text width=3.1cm, minimum height=0.82cm, align=center, font=\scriptsize},
  warncard/.style={draw=red!60!black, rounded corners=5pt, fill=red!5, text width=3.25cm, minimum height=0.82cm, align=center, font=\scriptsize}
]
\node[font=\bfseries\Large] at (0,4.78) {Distributed Data Table (DDT)};
\node[text width=11.9cm, align=center, font=\small] at (0,4.02) {The DDT maps content hashes to virtual holders; it stores neither content bytes nor physical node identities.};
\node[tablebox, minimum width=5.75cm, minimum height=2.75cm] (table) at (-3.85,1.35) {};
\node[font=\bfseries\small] at (-3.75,3.18) {Availability Records};
\draw (-6.50,2.18) rectangle (-1.20,0.42);
\draw (-4.62,2.18) -- (-4.62,0.42);
\draw (-6.50,1.58) -- (-1.20,1.58);
\draw (-6.50,1.00) -- (-1.20,1.00);
\node[font=\scriptsize\bfseries] at (-5.56,1.88) {Hash};
\node[font=\scriptsize\bfseries] at (-2.91,1.88) {Ephemeral VN holders};
\node[font=\scriptsize] at (-5.56,1.30) {$H(content_A)$};
\node[font=\scriptsize] at (-2.91,1.30) {$evn_1,\ evn_2$};
\node[font=\scriptsize] at (-5.56,0.72) {$H(content_B)$};
\node[font=\scriptsize] at (-2.91,0.72) {$evn_3,\ evn_4$};
\node[vnode] (a) at (0.10,2.20) {$evn_1$};
\node[vnode] (b) at (1.65,2.55) {$evn_2$};
\node[vnode] (c) at (3.20,2.05) {$evn_3$};
\node[vnode] (d) at (0.35,0.65) {$evn_4$};
\node[vnode] (e) at (1.85,0.25) {$evn_5$};
\node[vnode] (f) at (3.45,0.75) {$evn_6$};
\foreach \x/\y in {a/b,b/c,a/d,d/e,e/f,c/f,b/e} \draw[dashed, gray!65] (\x) -- (\y);
\coordinate (lookupSplit) at (-0.62,1.35);
\draw[thick, blue!60!black] (table.east) -- (lookupSplit);
\draw[arrow] (lookupSplit) -- (a.west);
\draw[arrow] (lookupSplit) -- (d.west);
\node[card] (interest) at (1.85,-1.15) {\textbf{Interest-driven}\\holders keep wanted content};
\draw[dashedarrow] (d.south) -- (interest.north west);
\draw[dashedarrow] (f.south) -- (interest.north east);
\node[warncard, text width=3.05cm] (noPhysical) at (6.10,1.35) {\textbf{No physical IDs}\\records expose only virtual holders};
\draw[dashedarrow] (c.east) -- (noPhysical.north west);
\draw[dashedarrow] (f.east) -- (noPhysical.south west);
\node[card] at (-4.55,-2.95) {\textbf{Hash check}\\bytes must match $H(content)$};
\node[card] at (-0.95,-2.95) {\textbf{Refresh}\\holders renew records};
\node[warncard] at (2.75,-2.95) {\textbf{No central delete}\\availability ends when holders stop};
\end{tikzpicture}}
\caption{Distributed Data Table.}
\end{figure}}

\newcommand{\thesisephemeraltikz}{%
\begin{figure}[H]
\centering
\resizebox{0.98\textwidth}{!}{%
\begin{tikzpicture}[
  arrow/.style={-{Latex[length=3mm]}, thick, blue!60!black},
  dashedarrow/.style={-{Latex[length=2.7mm]}, dashed, thick, gray!70},
  vnode/.style={circle, draw=gray!75!black, fill=gray!10, minimum size=0.95cm, align=center, font=\scriptsize},
  pnode/.style={rectangle, draw=gray!75!black, rounded corners=5pt, fill=gray!10, minimum width=1.85cm, minimum height=0.72cm, align=center, font=\scriptsize},
  data/.style={rectangle, draw=blue!60!black, rounded corners=4pt, fill=blue!6, minimum width=1.75cm, minimum height=0.7cm, align=center, font=\scriptsize},
  ddt/.style={rectangle, draw=green!55!black, rounded corners=5pt, fill=green!7, minimum width=2.5cm, minimum height=0.82cm, align=center, font=\scriptsize},
  card/.style={rectangle, draw=blue!55!black, rounded corners=5pt, fill=blue!5, text width=3.1cm, minimum height=0.82cm, align=center, font=\scriptsize},
  warncard/.style={rectangle, draw=red!60!black, rounded corners=5pt, fill=red!5, text width=3.1cm, minimum height=0.82cm, align=center, font=\scriptsize}
]
\node[font=\bfseries\Large] at (0,5.35) {Ephemeral Virtual Nodes};
\node[text width=10.6cm, align=center, font=\small] at (0,4.62) {Temporary holder identities avoid permanent publisher linkage.};
\node[font=\bfseries] at (-3.95,3.72) {Publication without publisher linkage};
\node[pnode] (publisher) at (-5.70,2.50) {Permanent VN\\publisher};
\node[data] (file) at (-3.35,2.50) {Encrypted\\content};
\node[ddt] (ddt) at (-0.88,2.50) {DDT record\\$H(content)$};
\node[vnode] (e1) at (1.45,3.18) {$evn_1$};
\node[vnode] (e2) at (3.00,2.64) {$evn_2$};
\node[vnode] (e3) at (1.45,1.72) {$evn_3$};
\node[vnode] (e4) at (4.45,2.50) {$evn_4$};
\draw[arrow] ([xshift=3pt]publisher.east) -- ([xshift=-4pt]file.west);
\draw[arrow] ([xshift=3pt]file.east) -- ([xshift=-4pt]ddt.west);
\draw[arrow, shorten >=2pt, shorten <=3pt] (ddt.north east) -- (e1.west);
\draw[arrow, shorten >=2pt, shorten <=3pt] (ddt.south east) -- (e3.west);
\foreach \a/\b in {e1/e2,e2/e4,e1/e3,e3/e2,e3/e4} \draw[dashed, gray!65] (\a) -- (\b);
\node[warncard] (pubcard) at (-4.75,0.78) {\textbf{Publisher not stored}\\DDT contains only holder references};
\node[card] (holdercard) at (2.95,0.50) {\textbf{Ephemeral holders}\\temporary VNs serve the content};
\draw (-6.35,-0.25) -- (5.90,-0.25);
\node[font=\bfseries] at (-3.95,-0.68) {Expiration and renewal};
\node[vnode] (old) at (-4.75,-1.90) {$evn_Q$};
\node[warncard] (expired) at (-1.75,-1.90) {\textbf{Expired}\\old record is ignored};
\node[vnode] (new) at (1.35,-1.90) {$evn_R$};
\node[ddt] (renewed) at (4.35,-1.90) {Renewed\\DDT record};
\draw[arrow, shorten >=2pt, shorten <=4pt] (old.east) -- (expired.west);
\draw[arrow, shorten >=2pt, shorten <=4pt] (expired.east) -- (new.west);
\draw[arrow, shorten >=2pt, shorten <=4pt] (new.east) -- (renewed.west);
\node[card] at (-3.95,-3.35) {\textbf{Refresh}\\holders renew records};
\node[card] at (0,-3.35) {\textbf{Persistence}\\content remains with holders};
\node[warncard] at (3.95,-3.35) {\textbf{Origin hidden}\\holders hide publisher};
\end{tikzpicture}}
\caption{Ephemeral Virtual Nodes.}
\end{figure}}

\newcommand{\thesisdtttikz}{%
\begin{figure}[H]
\centering
\resizebox{0.98\textwidth}{!}{%
\begin{tikzpicture}[
  block/.style={rectangle, rounded corners=5pt, minimum width=4.05cm, minimum height=4.10cm, text width=3.72cm, align=center},
  greenBlock/.style={block, draw=green!65!black, fill=green!10},
  blueBlock/.style={block, draw=blue!65!black, fill=blue!10},
  side/.style={rectangle, draw, rounded corners=5pt, minimum width=1.95cm, minimum height=3.45cm, text width=1.70cm, align=center},
  row/.style={rectangle, draw=gray!55, rounded corners=2pt, fill=white, minimum width=3.25cm, minimum height=0.36cm, align=center, font=\tiny},
  arrow/.style={-{Latex[length=3mm]}, thick, blue!55!black}
]
\node[font=\bfseries\Large, blue!35!black] at (0,5.75) {Content Discovery and Resolution Model (Two-Stage)};
\node[text width=13.5cm, align=center, blue!35!black] at (0,4.85) {A lightweight process connecting semantic discovery to current content availability using two distributed tables.};
\node[side] (query) at (-6.35,1.0) {\textbf{Query}\\[5pt]Tags\\[3pt]\texttt{tag1}\\\texttt{tag2}\\\texttt{tag3}};
\node[greenBlock] (dtt) at (-3.0,0.72) {};
\node[font=\bfseries\scriptsize, text width=3.3cm, align=center] at (-3.0,2.18) {1. Distributed Tag Table\\(DTT)};
\node[align=center, font=\scriptsize] at (-3.0,1.34) {Semantic Entry Point};
\node[row] at (-3.0,0.68) {\texttt{tag1} $\rightarrow$ \{hA, hB, hC\}};
\node[row] at (-3.0,0.08) {\texttt{tag2} $\rightarrow$ \{hB, hD, hE\}};
\node[row] at (-3.0,-0.52) {\texttt{tag3} $\rightarrow$ \{hA, hF, hG\}};
\node[blueBlock] (ddt) at (1.82,0.72) {};
\node[font=\bfseries\scriptsize, text width=3.3cm, align=center] at (1.82,2.18) {2. Distributed Data Table\\(DDT)};
\node[align=center, font=\scriptsize, text width=3.35cm] at (1.82,1.34) {Resolution and Availability Layer};
\node[row] at (1.82,0.68) {hA $\rightarrow$ descriptor + vn1};
\node[row] at (1.82,0.08) {hB $\rightarrow$ descriptor + vn2};
\node[row] at (1.82,-0.52) {hC $\rightarrow$ descriptor + vn3};
\node[side, minimum width=2.25cm, text width=1.95cm] (overlay) at (6.45,1.0) {\textbf{Overlay}\\\textbf{Network}\\[4pt]Storage VNs\\[3pt]vn1\\vn2\\vn3\\[4pt]Retrieve Object};
\draw[arrow] (query) -- (dtt);
\draw[arrow] (dtt.east) -- (ddt.west);
\node[font=\tiny, fill=white, inner sep=1.5pt] at (-0.59,1.12) {hashes};
\draw[arrow] (ddt.east) -- (overlay.west);
\node[font=\tiny, fill=white, inner sep=1.5pt] at (4.14,1.18) {holders};
\node[draw, rounded corners=4pt, fill=orange!10, minimum width=8.9cm, minimum height=0.72cm] at (0,-2.35) {\textbf{Overall Discovery Cycle:} \quad DTT $\rightarrow$ DDT $\rightarrow$ Overlay retrieval};
\node[draw=green!50!black, fill=green!8, rounded corners=4pt, text width=3.70cm, minimum height=1.65cm, align=left, font=\footnotesize] at (-4.45,-4.95) {\textbf{DTT}\\Maps tags to hashes.\\Semantic filtering.};
\node[draw=blue!50!black, fill=blue!8, rounded corners=4pt, text width=3.70cm, minimum height=1.65cm, align=left, font=\footnotesize] at (0,-4.95) {\textbf{DDT}\\Resolves hashes.\\Tracks holders.};
\node[draw=orange!70!black, fill=orange!8, rounded corners=4pt, text width=3.70cm, minimum height=1.65cm, align=left, font=\footnotesize] at (4.45,-4.95) {\textbf{Overlay Retrieval}\\Routes to storage VNs.\\Retrieves content.};
\end{tikzpicture}}
\caption{Content discovery and resolution model.}
\end{figure}}

\newcommand{\thesisdeliverytikz}{%
\begin{figure}[H]
\centering
\resizebox{0.98\textwidth}{!}{%
\begin{tikzpicture}[
  panel/.style={rectangle, draw, rounded corners=4pt, fill=gray!8, minimum width=2.2cm, minimum height=1.25cm, text width=2.15cm, align=center, font=\scriptsize},
  nodebox/.style={rectangle, draw, rounded corners=4pt, fill=blue!8, minimum width=2.05cm, minimum height=0.95cm, align=center},
  arrow/.style={-{Latex[length=2.8mm]}, thick},
  route/.style={-{Latex[length=3mm]}, very thick, blue!65!black},
  alt/.style={-{Latex[length=2.4mm]}, dashed, thick, blue!60!black},
  e2e/.style={-{Latex[length=3mm]}, very thick, green!55!black, bend left=20},
  reply/.style={-{Latex[length=2.8mm]}, dashed, thick, purple!60!black}
]
\node[font=\bfseries\Large, blue!35!black] at (0,5.05) {Message Routing Process};
\node[font=\bfseries] at (0,4.58) {Secure, Decentralized, and Obfuscated Communication};
\foreach \i/\x/\t in {1/-5.0/DRT Query,2/-2.5/Entry Points,3/0/Hop-by-Hop Routing,4/2.5/Protected Channel,5/5.0/Return Path} {
  \node[panel] at (\x,3.55) {\textbf{\i. \t}\\\scriptsize workflow step};
}

\node[nodebox] (sender) at (-5.90,1.0) {Sender\\Virtual A};
\node[cylinder, shape border rotate=90, draw, aspect=0.25, minimum height=1.1cm, minimum width=0.7cm, fill=gray!20] (drt) at (-3.65,1.0) {DRT};
\node[ellipse, draw, minimum width=6.25cm, minimum height=3.45cm, fill=blue!5] (cloud) at (0.15,0.05) {};
\foreach \n/\x/\y/\label in {a/-1.65/0.30/A,b/-0.35/1.20/B,c/1.0/0.90/C,d/-0.15/-0.40/D,e/-1.35/-0.95/E,f/1.45/-0.85/F,g/0.45/-1.35/G,h/2.45/0.10/H} {
  \node[rectangle, draw, rounded corners=2pt, fill=gray!25,
        minimum width=0.82cm, minimum height=0.58cm, font=\tiny] (n\n) at (\x,\y) {PN \label};
}
\node[nodebox, draw=purple!70!black] (dest) at (5.55,1.0) {Destination\\Virtual B};
\draw[arrow] (sender) -- (drt);
\node[font=\tiny, fill=white, inner sep=1pt, align=center] at (-4.48,1.34) {DRT\\query};
\draw[arrow, dashed] (drt) -- (na);
\node[font=\scriptsize, fill=white, inner sep=1pt] at (-2.54,1.33) {entry points};
\draw[route] (na) -- (nb) -- (nc) -- (nh) -- (dest);
\draw[alt] (na) -- (ne);
\draw[alt] (ne) -- (nd);
\draw[alt] (nd) -- (nf);
\draw[alt] (nf) -- (nh);
\draw[alt] (nd) -- (nc);
\draw[e2e] (sender.north) to node[above, font=\scriptsize, fill=white, inner sep=1pt] {End-to-end encrypted session} (dest.north);
\draw[reply] (dest.south) .. controls (4.0,-2.92) and (-3.0,-2.92) .. (sender.south);
\node[draw, rounded corners=4pt, fill=blue!7, text width=2.70cm, minimum height=1.8cm, align=left, font=\footnotesize] at (-4.8,-3.85) {\textbf{Local knowledge}\\Previous hop and next hop only.};
\node[draw, rounded corners=4pt, fill=blue!7, text width=2.70cm, minimum height=1.8cm, align=left, font=\footnotesize] at (-1.6,-3.85) {\textbf{Protection}\\E2E layer plus hop layer.};
\node[draw, rounded corners=4pt, fill=blue!7, text width=2.70cm, minimum height=1.8cm, align=left, font=\footnotesize] at (1.6,-3.85) {\textbf{Failure}\\Retry entry point or route.};
\node[draw, rounded corners=4pt, fill=blue!7, text width=2.70cm, minimum height=1.8cm, align=left, font=\footnotesize] at (4.8,-3.85) {\textbf{Properties}\\Obfuscated path, no global route.};
\end{tikzpicture}}
\caption{Message routing process.}
\end{figure}}

\newcommand{\thesisvirtualcommunicationflowtikz}{%
\begin{figure}[H]
\centering
\resizebox{\textwidth}{!}{%
\begin{tikzpicture}[
    node distance=7mm,
    stage/.style={draw, rounded corners=2pt, align=center, minimum width=2.65cm,
        minimum height=1.05cm, fill=blue!5, font=\small},
    flow/.style={-{Latex[length=2.2mm]}, thick, draw=blue!55!black},
    note/.style={align=center, font=\scriptsize, text width=2.85cm}
]
\node[stage] (source) {Virtual node A\\logical origin};
\node[stage, right=of source] (lookup) {DPNT / DRT\\resolution};
\node[stage, right=of lookup] (route) {Route strategy and\\hop-by-hop route};
\node[stage, right=of route] (session) {End-to-end\\virtual session};
\node[stage, right=of session] (target) {Virtual node B\\logical destination};
\draw[flow] (source) -- (lookup);
\draw[flow] (lookup) -- (route);
\draw[flow] (route) -- (session);
\draw[flow] (session) -- (target);
\node[note, below=5mm of lookup] {Discovers physical entry points and published routes.};
\node[note, below=5mm of route] {Selects hops, constructs local mappings, and carries \texttt{ROUTE\_DATA}.};
\node[note, below=5mm of session] {Establishes keys and carries protected messages.};
\end{tikzpicture}}
\caption{From distributed discovery to virtual application communication.}
\label{fig:virtual-communication-flow}
\end{figure}}

\newcommand{\thesiscryptographicflowtikz}{%
\begin{figure}[H]
\centering
\resizebox{\textwidth}{!}{%
\begin{tikzpicture}[
    node distance=8mm and 10mm,
    stage/.style={draw, rounded corners=2pt, align=center, text width=3.15cm,
        minimum height=1.1cm, fill=green!5, font=\small},
    flow/.style={-{Latex[length=2.2mm]}, thick, draw=green!40!black},
    support/.style={draw, rounded corners=2pt, align=center, text width=3.15cm,
        minimum height=1.05cm, fill=orange!8, font=\small}
]
\node[stage] (identity) {Identity\\SHA-512 identifier};
\node[stage, right=of identity] (auth) {Authentication\\ML-DSA-65 signatures};
\node[stage, right=of auth] (kem) {Key establishment\\ML-KEM-768};
\node[stage, right=of kem] (aead) {Protected session\\AES-256-GCM-SIV + AAD};
\draw[flow] (identity) -- (auth);
\draw[flow] (auth) -- (kem);
\draw[flow] (kem) -- (aead);
\node[support, below=9mm of auth] (records) {Signed control-plane and distributed records};
\node[support, below=9mm of kem] (pow) {Publication proof-of-work\\canonical payload + nonce};
\draw[flow] (auth) -- (records);
\draw[flow] (pow) -- (records);
\coordinate (powroute) at ($(pow.south)+(0,-5mm)$);
\draw[flow] (identity.south) |- (powroute) -- (pow.south);
\end{tikzpicture}}
\caption{Security mechanisms and the stages at which they are applied.}
\label{fig:cryptographic-flow}
\end{figure}}
